\section{Experimental Methodology}
\subsection{Data sets} 
We run whole genome tests on mrFAST with one real and two simulated data sets. We obtained the biological data sets from the 1000 Genomes Project Phase I~\cite{10002012integrated} and produced simulated reads using Mason~\cite{holtgrewe2010mason} at different read lengths. In all of our tests, we use the reference GRCh37 produced by the 1000 Genomes Project~\cite{10002015global}. In order to have a fair comparison for accuracy, we use the same data sets from GateKeeper~\cite{Alser2017} for the accuracy tests. 

For filtering throughput and accuracy analyses, the datasets contain 30 million read and candidate reference segment pairs, seeded by mrFAST, using biological short read sets~(downloaded from European Nucleotide Archive) with different error threshold values. Additionally, we generate read and candidate reference segment pairs using two of the state-of-the-art mapping tools, Minimap2~\cite{li2018minimap2}, and BWA-MEM~\cite{li2013aligning} for accuracy analyses. For Minimap2, we extract the pairs just before the first chaining function ($mm\_chain\_dp$) since this is the first dynamic programming part and gather the samples in 11 different sets for error threshold from 0 to 10, each containing 30 million pairs. For BWA-MEM, we extract the pairs before the final global alignment call (\textit{ksw\_global2}). BWA-MEM generates much fewer pairs than 30 million at this point, therefore the data set size for each error threshold~(\textit{e} = {0, \ldots, 10}) is different.  
By using Edlib's global alignment mode, we prepare the filtering status of data sets to maintain consistency on ground truth. In all of the experiments, the maximum error threshold is $10\%$ of the read length. Please refer to Sup. Table S.1 for the details about the data sets.
%%%%%%%%%%%%%%%%%%%%%%%%%%%%%%%%%%%%%%%%%%%%%%%%%%%%%%%%%%
%%%%%%%%%%%%%%%%%%%%%%%%%%%%%%%%%%%%%%%%%%%%%%%%%%%%%%%%%%
%\subsection{Device Specifications}
\subsection{Experimental Setup}
We run our performance experiments in two different setups. \textit{Setup\_1} includes a 2.30GHz Intel Xeon Gold 6140 CPU with 754G RAM. There are 8 NVIDIA GeForce GTX 1080 Ti GPUs~(Pascal architecture), each with 10GB global memory connected to this processor. CUDA compute capability of the devices is 6.1 with CUDA driver v10.1 installed. Data transfer between the host and devices is managed by PCIe generation 3 with 16 lanes. In \textit{Setup\_2}, we use a 3.30GHz Intel Xeon CPU E5-2643 0 processor with 256G RAM. 4 NVIDIA Tesla K20X GPGPUs are connected to this host. Each device has 5GB global memory and the data transfer is maintained via PCIe generation 2 with 16 lanes. CUDA compute capability of the devices is 3.5 with CUDA driver v10.2. Tesla K20X has Kepler architecture, therefore data prefetching is not supported in Setup\_2. In all of our tests, we enable persistence mode in GPU to keep the devices initialized.

%%%%%%%%%%%%%%%%%%%%%%%%%%%%%%%%%%%%%%%%%%%%%%%%%%%%%%%%%%
%%%%%%%%%%%%%%%%%%%%%%%%%%%%%%%%%%%%%%%%%%%%%%%%%%%%%%%%%%
\subsection{Filtering Throughput Evaluation}
We run our throughput analysis on the datasets containing read and reference segment pairs. Each dataset collected for this purpose includes 30 million pairs in total. We calculate the runtime taken for the filtration of a single pair out of 30 million, then determine the total number of pairs that can be filtered in 40 minutes to have a fair throughput comparison with the other tools. 

We report two different time measurements for throughput analysis: \textit{kernel time} and \textit{filter time}. \textit{Kernel time} denotes the time taken only by GPU devices and we record this time using CUDA Event API. Since GateKeeper-GPU uses batched kernel calls, we add all kernel times in execution and report the sum. \textit{Filter time} represents the total time spent for filtering, including host operations such as data transfer and encoding the sequences. Therefore, we measure filter time from the host's perspective. For both of these measurements, we also provide different results for encoding the pairs by the host~(CPU) and the device~(GPU). In multi-GPU throughput analysis, kernel time represents the time of the device, which takes the longest time to complete among all other active devices. In all of our timing experiments, we run the tests ten times and report the arithmetic mean to minimize the effect of random experimental errors on results. 

We compare GateKeeper-GPU's filtering throughput with its CPU version comprehensively and make a brief comparison with its FPGA version~\cite{Alser2017}. To maintain fairness as much as possible, we implement GateKeeper-CPU in a multicore fashion and report the results of 12 cores.   

%%%%%%%%%%%%%%%%%%%%%%%%%%%%%%%%%%%%%%%%%%%%%%%%%%%%%%%%%%
%%%%%%%%%%%%%%%%%%%%%%%%%%%%%%%%%%%%%%%%%%%%%%%%%%%%%%%%%%
\subsection{Accuracy Evaluation} 
In our accuracy analyses, we consider Edlib's~\cite{vsovsic2017edlib} edit distance as the ground truth and calculate the edit distance between the read and reference segments by using Edlib's global alignment. According to the edit distance, we produce a filtering status for each pair as \textit{reject} if the edit distance is larger than the threshold or \textit{accept} if otherwise. Throughout accuracy experiments, we analyze \textit{false accept}, \textit{false reject}, and \textit{true reject} counts. A false accept represents a read and reference segment pair that Edlib rejects because of exceeding the error threshold, but is accepted by the filter. On the contrary, a false reject case is a valid pair with fewer errors than the threshold but is rejected by the filter. True rejects are the pairs that are rejected by both Edlib and GateKeeper-GPU.

We have two approaches for testing the accuracy of GateKeeper-GPU. First, we record the filtering status for the datasets described in Section 4.1 and compare the results with Edlib. Since GateKeeper-GPU gives a direct pass to the pairs that contain unknown base call characters~(\textit{`N'}), to be able to observe the actual accept and reject counts, we exclude these pairs from the tests and report the comparison with Edlib accordingly. For the sake of simplicity, we will call these pairs \textit{undefined} for the rest of the work. 

Second, we compare GateKeeper-GPU with the original GateKeeper FPGA implementation, SHD, MAGNET, Shouji, and SneakySnake in terms of false accept and false reject counts. We denote the original GateKeeper implementation as `GateKeeper-FPGA' for labeling. We choose the highest-edit and lowest-edit profile datasets of three different read lengths for these tests. For 100bp, the lowest-edit containing data set is prepared by mrFAST's seeding error threshold \textit{e =}~2, and the highest-edit data set is curated by error threshold \textit{e =}~40. Likewise, mrFAST error thresholds for the lowest-edit and the highest-edit profile data sets are \textit{e =}~4 and \textit{e =}~70 for 150bp; \textit{e =}~8 and \textit{e =}~100 for 250bp, respectively. Because the other tools do not have distinguishing mechanisms for undefined pairs, to maintain a fair comparison in this test series, we include these pairs in GateKeeper-GPU's results and report the numbers accordingly. 

%%%%%%%%%%%%%%%%%%%%%%%%%%%%%%%%%%%%%%%%%%%%%%%%%%%%%%%%%%
%%%%%%%%%%%%%%%%%%%%%%%%%%%%%%%%%%%%%%%%%%%%%%%%%%%%%%%%%%
\subsection{Whole Genome Evaluation} 
We evaluate GateKeeper-GPU's performance and accuracy in the full workflow of a read mapping tool by integrating GateKeeper-GPU into mrFAST. GateKeeper was also integrated into mrFAST for testing. Therefore, testing GateKeeper-GPU as a part of mrFAST enables us to briefly compare it with the original work. 

We use the same datasets that GateKeeper~\cite{Alser2017} uses in whole genome comparison experiments: one real 100bp data set and one simulated 300bp data set~(i.e., sim\_set\_1). We add another simulated 150bp data set~(i.e., sim\_set\_2) analysis to GateKeeper-GPU's results. We collect the following metrics from alignment for evaluation: the number of mappings, the number of mapped reads, the total number of candidate mappings, the total number of candidate mappings that enter verification, time spent for verification, time spent for preprocessing before pre-alignment filtering, and total kernel time spent for running GateKeeper-GPU. We use a single GPU in these experiments and provide results for encoding in both device and host design. 
%%%%%%%%%%%%%%%%%%%%%%%%%%%%%%%%%%%%%%%%%%%%%%%%%%%%%%%%%%
%%%%%%%%%%%%%%%%%%%%%%%%%%%%%%%%%%%%%%%%%%%%%%%%%%%%%%%%%%
\subsection{Resource Utilization and Power Evaluation} 
Warp occupancy is one of the important indicators when evaluating the kernel performance of a CUDA application. A \textit{warp} is a group of 32 adjacent threads in a block, and Streaming Multiprocessor~(i.e., SM) schedules the same instruction for the threads in a warp~\cite{achieved_occupancy}. Warp occupancy is the ratio of active warps over the maximum supported number of warps for SM. The number of warps, blocks, registers per SM, and shared memory determines theoretical occupancy. With these conditions, the maximum number of registers per thread is 32 for $100\%$ occupancy while using all threads in a warp. GateKeeper-GPU does not utilize shared memory, and we opt to maximize the number of warps and blocks to maintain a high throughput. Using the CUDA occupancy calculator~\cite{cudaoccupancycalculator}, we calculate GateKeeper-GPU's theoretical warp occupancy and present a comparison with the achieved occupancy value. 

We perform profiling experiments on the 100bp and 250bp data sets with error thresholds \textit{e = }4 and \textit{e = }10, respectively, by using CUDA command-line profiler nvprof \cite{nvprof}. By conducting system profiling and metrics analyses, we provide GateKeeper-GPU's power consumption, warp execution efficiency, and multiprocessor activity.